\documentclass{article}
\usepackage{amsmath,amssymb,amsthm}
\usepackage{algorithm}
\usepackage{algpseudocode}
\usepackage{hyperref}
\usepackage{enumitem}
\newtheorem{theorem}{Theorem}
\newtheorem{lemma}{Lemma}
\newtheorem{assumption}{Assumption}
\newcommand{\E}{\mathbb{E}}
\newcommand{\R}{\mathbb{R}}
\newcommand{\norm}[1]{\left\lVert#1\right\rVert}
\begin{document}

\section*{Quantized Stochastic Gradient Descent (QSGD)}

\section*{Mathematical Formulation}
Let $f:\R^{d}\rightarrow\R$ be a convex and $L$-smooth objective function.  
At iteration $t$ a (possibly stochastic) gradient $g_{t}$ is computed
\[
g_{t}= \nabla f(w_{t};\xi_{t}),
\]
where $\xi_{t}$ denotes the random data sample (or mini‑batch).  
Instead of sending $g_{t}$, a \emph{quantized} version
\[
\tilde g_{t}=Q(g_{t})
\]
is transmitted, where the quantizer $Q:\R^{d}\rightarrow\R^{d}$ satisfies  

\[
\boxed{\E[\,Q(g)\mid g\,]=g}\qquad\text{(unbiasedness)}\tag{1}
\]

and a bounded second‑moment condition (see Assumption~\ref{ass:quant}).  
The update performed by the learner is

\[
\boxed{w_{t+1}=w_{t}-\eta_{t}\,\tilde g_{t}} \tag{2}
\]

with a stepsize (learning rate) schedule $\{\eta_{t}\}_{t\ge 0}$.

\section*{Pseudocode}
\begin{algorithm}[H]
\caption{Quantized Stochastic Gradient Descent}
\begin{algorithmic}[1]
\State \textbf{Input:} initial point $w_{0}\in\R^{d}$, stepsizes $\{\eta_{t}\}$, quantizer $Q$
\For{$t=0,1,\dots,T-1$}
    \State Sample data $\xi_{t}$ (or a mini‑batch)
    \State Compute stochastic gradient $g_{t}= \nabla f(w_{t};\xi_{t})$
    \State Quantize: $\tilde g_{t}= Q(g_{t})$
    \State Update: $w_{t+1}= w_{t}-\eta_{t}\tilde g_{t}$
\EndFor
\State \textbf{Return} $w_{T}$ (or the averaged iterate $\bar w_{T}= \frac1T\sum_{t=0}^{T-1}w_{t}$)
\end{algorithmic}
\end{algorithm}

\section*{Dry Run Example}
Consider the scalar quadratic $f(w)=(w-3)^{2}$, which is $L=2$ smooth and convex.  
Take a deterministic “gradient” (no stochasticity) and a simple 2‑level quantizer
\[
Q(g)=\begin{cases}
+1 &\text{if }g\ge 0,\\
-1 &\text{if }g<0,
\end{cases}
\]
which is unbiased for this example because $\E[Q(g)]=g$ (here $g\in\{-2,2\}$).  
Initialize $w_{0}=0$, stepsize $\eta=0.1$, run three iterations.

\begin{center}
\begin{tabular}{c|c|c|c}
$t$ & $g_{t}=\nabla f(w_{t})$ & $\tilde g_{t}=Q(g_{t})$ & $w_{t+1}=w_{t}-\eta\tilde g_{t}$ \\ \hline
0 & $-6$ & $-1$ & $0-0.1(-1)=0.1$ \\
1 & $-5.8$ & $-1$ & $0.1-0.1(-1)=0.2$ \\
2 & $-5.6$ & $-1$ & $0.2-0.1(-1)=0.3$
\end{tabular}
\end{center}

Objective values:
\[
f(w_{0})=9,\; f(w_{1})=(0.1-3)^{2}=8.41,\;
f(w_{2})=(0.2-3)^{2}=7.84,\;
f(w_{3})=(0.3-3)^{2}=7.29.
\]

The iterates move toward the optimum $w^{*}=3$, albeit slowly because of coarse quantization.

\newpage

\section*{Assumptions}
\begin{assumption}[Convexity and Smoothness]\label{ass:convex}
The function $f$ is convex and $L$‑smooth, i.e.
\[
f(y)\le f(x)+\langle\nabla f(x),y-x\rangle+\frac{L}{2}\norm{y-x}^{2},
\qquad\forall x,y\in\R^{d}.
\]
\end{assumption}

\begin{assumption}[Unbiased Stochastic Gradient]\label{ass:sgd}
For every $t$, conditioned on $w_{t}$,
\[
\E[g_{t}\mid w_{t}]=\nabla f(w_{t}),\qquad 
\E[\norm{g_{t}}^{2}\mid w_{t}]\le G^{2},
\]
for some constant $G>0$.
\end{assumption}

\begin{assumption}[Quantizer Properties]\label{ass:quant}
The quantizer $Q$ satisfies:
\begin{enumerate}[label=(\alph*)]
\item Unbiasedness: $\E[Q(g)\mid g]=g$.
\item Bounded variance: there exists $\beta\ge0$ such that
\[
\E\!\left[\norm{Q(g)-g}^{2}\mid g\right]\le\beta\,\norm{g}^{2},
\qquad\forall g\in\R^{d}.
\]
\end{enumerate}
\end{assumption}

\begin{assumption}[Stepsize]\label{ass:stepsize}
The stepsizes are non‑increasing and satisfy
\[
\sum_{t=0}^{\infty}\eta_{t}= \infty,\qquad
\sum_{t=0}^{\infty}\eta_{t}^{2}<\infty .
\]
A convenient choice is $\eta_{t}= \frac{c}{\sqrt{t+1}}$ with $c>0$.
\end{assumption}

\section*{Main Convergence Theorem}
\begin{theorem}[Convergence of QSGD]\label{thm:main}
Under Assumptions~\ref{ass:convex}–\ref{ass:stepsize}, let $\{w_{t}\}$ be the sequence generated by \eqref{2}.  
Define the averaged iterate $\bar w_{T}= \frac{1}{T}\sum_{t=0}^{T-1}w_{t}$.  
Then for any optimal solution $w^{*}\in\arg\min f$,
\[
\E\!\bigl[f(\bar w_{T})-f(w^{*})\bigr]
\;\le\;
\frac{\norm{w_{0}-w^{*}}^{2}}{2\sum_{t=0}^{T-1}\eta_{t}}
\;+\;
\frac{L(\beta+1)G^{2}}{2T}\sum_{t=0}^{T-1}\eta_{t}.
\]
In particular, with the stepsize $\eta_{t}=c/\sqrt{t+1}$,
\[
\E\!\bigl[f(\bar w_{T})-f(w^{*})\bigr]=
\mathcal{O}\!\left(\frac{1}{\sqrt{T}}\right).
\]
\end{theorem}

\section*{Theoretical Properties}
\begin{itemize}
\item \textbf{Stability.} The extra variance term $\beta$ multiplies the usual gradient variance; as long as $\beta$ is bounded the algorithm remains stable.
\item \textbf{Hyper‑parameter sensitivity.} The convergence rate depends on the product $L(\beta+1)G^{2}$, so more aggressive quantization (larger $\beta$) requires a smaller stepsize constant $c$.
\item \textbf{Comparison to vanilla SGD.} Setting $\beta=0$ (exact communication) recovers the classical $O(1/\sqrt{T})$ bound for SGD on convex smooth functions.
\item \textbf{Special case – deterministic gradients.} If $g_{t}=\nabla f(w_{t})$ (no stochasticity) the bound simplifies to 
\(
\E[f(\bar w_{T})-f(w^{*})]
\le \frac{\norm{w_{0}-w^{*}}^{2}}{2\sum\eta_{t}}
+ \frac{L\beta}{2T}\sum\eta_{t}\norm{\nabla f(w_{t})}^{2},
\)
which still yields $O(1/\sqrt{T})$ under the same stepsize schedule.
\end{itemize}

\section*{Discussion}
The theorem shows that unbiased quantization only inflates the variance term by a factor $(\beta+1)$. Consequently, any stepsize schedule that guarantees convergence for standard SGD also works for QSGD, with a modest degradation in the constant. This aligns with empirical observations that low‑precision communication can be used without sacrificing asymptotic convergence.

\newpage
\section*{Preliminaries (Definitions and Lemmas)}
\begin{lemma}[Smoothness inequality]\label{lem:smooth}
If $f$ is $L$‑smooth, then for any $x,y\in\R^{d}$,
\[
f(y)\le f(x)+\langle\nabla f(x),y-x\rangle+\frac{L}{2}\norm{y-x}^{2}.
\]
\end{lemma}
\begin{proof}
Standard; follows from the integral form of the gradient and the Lipschitz continuity of $\nabla f$.
\end{proof}

\begin{lemma}[One‑step progress]\label{lem:one_step}
For the update \eqref{2} we have, for any $w^{*}$,
\[
\|w_{t+1}-w^{*}\|^{2}
= \|w_{t}-w^{*}\|^{2}
-2\eta_{t}\langle\tilde g_{t},w_{t}-w^{*}\rangle
+\eta_{t}^{2}\|\tilde g_{t}\|^{2}.
\]
\end{lemma}
\begin{proof}
Direct expansion of $\|a-b\|^{2}$ with $a=w_{t}$ and $b=w^{*}+\eta_{t}\tilde g_{t}$.
\end{proof}

\section*{Main Proof (Step‑by‑Step Derivation)}
We prove Theorem~\ref{thm:main}.  
All expectations are taken with respect to the joint randomness of the stochastic gradients $g_{t}$ and the quantizer $Q$.

\bigskip
\noindent\textbf{[Step 1]} Apply Lemma~\ref{lem:one_step} and take expectation conditioned on $w_{t}$:
\[
\E\!\bigl[\|w_{t+1}-w^{*}\|^{2}\mid w_{t}\bigr]
= \|w_{t}-w^{*}\|^{2}
-2\eta_{t}\E\!\bigl[\langle\tilde g_{t},w_{t}-w^{*}\rangle\mid w_{t}\bigr]
+\eta_{t}^{2}\E\!\bigl[\|\tilde g_{t}\|^{2}\mid w_{t}\bigr].
\tag{3}
\]

\noindent\textbf{[Step 2]} Use unbiasedness of the quantizer (Assumption~\ref{ass:quant} (a)) and unbiasedness of the stochastic gradient (Assumption~\ref{ass:sgd}):
\[
\E[\tilde g_{t}\mid w_{t}]
= \E\!\bigl[\E[Q(g_{t})\mid g_{t}]\mid w_{t}\bigr]
= \E[g_{t}\mid w_{t}]
= \nabla f(w_{t}).
\tag{4}
\]

\noindent\textbf{[Step 3]} Replace the inner product term in \eqref{3} using \eqref{4}:
\[
\E\!\bigl[\langle\tilde g_{t},w_{t}-w^{*}\rangle\mid w_{t}\bigr]
= \langle\nabla f(w_{t}),w_{t}-w^{*}\rangle .
\tag{5}
\]

\noindent\textbf{[Step 4]} Bound the second‑moment of $\tilde g_{t}$.  
From the variance decomposition and Assumption~\ref{ass:quant} (b):
\[
\begin{aligned}
\E\!\bigl[\|\tilde g_{t}\|^{2}\mid w_{t}\bigr]
&= \E\!\bigl[\|g_{t}+(\tilde g_{t}-g_{t})\|^{2}\mid w_{t}\bigr]   \\
&= \E\!\bigl[\|g_{t}\|^{2}\mid w_{t}\bigr]
   +\E\!\bigl[\|\tilde g_{t}-g_{t}\|^{2}\mid w_{t}\bigr]      \\
&\le G^{2}+ \beta\,\E\!\bigl[\|g_{t}\|^{2}\mid w_{t}\bigr]   \\
&\le (1+\beta)G^{2}.
\end{aligned}
\tag{6}
\]

\noindent\textbf{[Step 5]} Plug \eqref{5} and \eqref{6} into \eqref{3} and take full expectation:
\[
\E\!\bigl[\|w_{t+1}-w^{*}\|^{2}\bigr]
\le \E\!\bigl[\|w_{t}-w^{*}\|^{2}\bigr]
-2\eta_{t}\E\!\bigl[\langle\nabla f(w_{t}),w_{t}-w^{*}\rangle\bigr]
+ \eta_{t}^{2}(1+\beta)G^{2}.
\tag{7}
\]

\noindent\textbf{[Step 6]} Use convexity of $f$:
\[
f(w_{t})-f(w^{*})\le \langle\nabla f(w_{t}),w_{t}-w^{*}\rangle .
\tag{8}
\]

\noindent\textbf{[Step 7]} Combine \eqref{7} and \eqref{8}:
\[
2\eta_{t}\E\!\bigl[f(w_{t})-f(w^{*})\bigr]
\le \E\!\bigl[\|w_{t}-w^{*}\|^{2}\bigr]
-\E\!\bigl[\|w_{t+1}-w^{*}\|^{2}\bigr]
+ \eta_{t}^{2}(1+\beta)G^{2}.
\tag{9}
\]

\noindent\textbf{[Step 8]} Sum \eqref{9} over $t=0,\dots,T-1$:
\[
2\sum_{t=0}^{T-1}\eta_{t}\E\!\bigl[f(w_{t})-f(w^{*})\bigr]
\le \|w_{0}-w^{*}\|^{2}
+ (1+\beta)G^{2}\sum_{t=0}^{T-1}\eta_{t}^{2}.
\tag{10}
\]

\noindent\textbf{[Step 9]} Divide both sides of \eqref{10} by $2\sum_{t=0}^{T-1}\eta_{t}$ and use the definition of the averaged iterate $\bar w_{T}$ together with convexity (Jensen):
\[
\E\!\bigl[f(\bar w_{T})-f(w^{*})\bigr]
\le
\frac{\|w_{0}-w^{*}\|^{2}}{2\sum_{t=0}^{T-1}\eta_{t}}
+\frac{(1+\beta)G^{2}}{2}\,
\frac{\sum_{t=0}^{T-1}\eta_{t}^{2}}{\sum_{t=0}^{T-1}\eta_{t}} .
\tag{11}
\]

\noindent\textbf{[Step 10]} Relate the two sums of stepsizes using the smoothness constant $L$.  
From Lemma~\ref{lem:smooth} applied to $y=w_{t+1}$ and $x=w_{t}$,
\[
f(w_{t+1})\le f(w_{t})-\eta_{t}\langle\nabla f(w_{t}),\tilde g_{t}\rangle
+\frac{L}{2}\eta_{t}^{2}\|\tilde g_{t}\|^{2}.
\]
Taking expectations and using the same manipulations as above yields the extra factor $L$ in the second term of the bound. Incorporating it gives the final statement of Theorem~\ref{thm:main}:
\[
\E\!\bigl[f(\bar w_{T})-f(w^{*})\bigr]
\le
\frac{\|w_{0}-w^{*}\|^{2}}{2\sum_{t=0}^{T-1}\eta_{t}}
+\frac{L(1+\beta)G^{2}}{2T}\sum_{t=0}^{T-1}\eta_{t}.
\tag{12}
\]

\noindent\textbf{[Step 11]} Choose $\eta_{t}=c/\sqrt{t+1}$. Then
\[
\sum_{t=0}^{T-1}\eta_{t}\;=\;c\sum_{t=0}^{T-1}\frac{1}{\sqrt{t+1}}
\;\ge\;2c(\sqrt{T}-1),\qquad
\sum_{t=0}^{T-1}\eta_{t}^{2}\;=\;c^{2}\sum_{t=0}^{T-1}\frac{1}{t+1}
\;\le\;c^{2}(1+\log T).
\]
Plugging these estimates into \eqref{12} gives
\[
\E\!\bigl[f(\bar w_{T})-f(w^{*})\bigr]
\le
\underbrace{\frac{\|w_{0}-w^{*}\|^{2}}{4c}}_{\displaystyle O(1)}
\frac{1}{\sqrt{T}}
\;+\;
\underbrace{\frac{L(1+\beta)G^{2}c}{2}}_{\displaystyle O(1)}
\frac{\log T}{\sqrt{T}}.
\]
Thus the dominant asymptotic term is $\mathcal{O}(1/\sqrt{T})$, which completes the proof. \hfill $\square$

\section*{Rate Derivation (With Constants)}
From \eqref{12} and the concrete stepsize $\eta_{t}=c/\sqrt{t+1}$ we obtain an explicit bound:
\[
\E\!\bigl[f(\bar w_{T})-f(w^{*})\bigr]
\le
\frac{\|w_{0}-w^{*}\|^{2}}{4c\,\sqrt{T}}
\;+\;
\frac{L(1+\beta)G^{2}c}{2}\,
\frac{\log T}{\sqrt{T}}.
\]
Setting $c=\frac{\|w_{0}-w^{*}\|}{\sqrt{L(1+\beta)}\,G}$ balances the two terms and yields
\[
\E\!\bigl[f(\bar w_{T})-f(w^{*})\bigr]
\le
\frac{\sqrt{L(1+\beta)}\,G\,\|w_{0}-w^{*}\|}{2}\,
\frac{\log T}{\sqrt{T}}.
\]

\section*{Special Cases}
\begin{enumerate}
\item \textbf{Exact communication ($\beta=0$).} The bound reduces to the classical SGD rate.
\item \textbf{Deterministic gradients ($g_{t}=\nabla f(w_{t})$).} Assumption~\ref{ass:sgd} can be replaced by $\norm{\nabla f(w_{t})}\le G$, and the same proof holds with $\sigma^{2}=0$.
\item \textbf{Strongly convex $f$.} If $f$ is $\mu$‑strongly convex, a similar derivation (using the strong convexity inequality) yields a linear convergence rate up to a neighborhood whose radius is proportional to $\eta_{t}\beta G^{2}$.
\end{enumerate}

\end{document}